二级市场部分不应以“没有数据”回避判断。以下价格、成交、融资融券、沪股通、基金和龙虎榜数据均已保留对应的原始结构化响应；各指标的日期不同，不能把它们伪装成同一时点的资金流，也不能把筹码变化直接等同于基本面变化。结论是：截至2026年7月22日，价格处于中短期走弱后的反弹窗口，融资余额同步回落；季度沪股通与基金持仓仍是正持仓，但二者均是披露/结算时点的存量，而非当日净买入。

\section{复权价格、趋势与回撤：先把口径说清}

连续历史数据来自东方财富公开日线端点，主序列为\texttt{fqt=1}的\textbf{前复权}日K，覆盖2023年7月3日至2026年7月22日共727个交易日。报告中的收益、均线、区间高低和最大回撤均只使用这条前复权序列；不复权（\texttt{fqt=0}）和后复权（\texttt{fqt=2}）原始响应已同时归档，用于核验而不混算。\srcid{PR-02}

\begin{exhibitbox}[连续前复权日K：价格表现与回撤]
\small
\begin{tabularx}{\exhibitboxwidth}{L{2.1cm}L{3.0cm}R{2.2cm}R{1.6cm}L{3.2cm}X}
\toprule
\textbf{区间} & \textbf{实际起止交易日} & \textbf{起/末收盘（元）} & \textbf{收益} & \textbf{区间高/低（元）} & \textbf{收盘最大回撤} \\
\midrule
1个月 & 2026-06-22--07-22 & 36.97 / 33.02 & -10.68\% & 37.90 / 31.76 & -14.22\% \\
3个月 & 2026-04-22--07-22 & 36.12 / 33.02 & -8.58\% & 43.06 / 31.76 & -22.77\% \\
年初至今 & 2026-01-05--07-22 & 34.12 / 33.02 & -3.22\% & 43.06 / 29.64 & -22.77\% \\
1年 & 2025-07-22--2026-07-22 & 34.34 / 33.02 & -3.84\% & 43.06 / 29.64 & -22.89\% \\
3年覆盖 & 2023-07-03--2026-07-22 & 31.90 / 33.02 & +3.51\% & 45.34 / 23.99 & -37.43\% \\
\bottomrule
\end{tabularx}
\sourcenote{PR-02；最大回撤为区间内前复权每日收盘价相对此前收盘峰值的最大跌幅，不含税费、基准和总回报。}
\end{exhibitbox}

截止日33.02元仅略高于MA5的32.72元，却低于MA10/20/60/120/250的33.72/34.37/36.55/35.38/35.16元，分别低2.08\%、3.92\%、9.65\%、6.66\%和6.10\%。近5/20/60个交易日区间依次为31.76--33.80元、31.76--37.90元和31.76--43.06元。它支持的最窄结论是：短线在7月17日低点后出现反弹，但中期仍位于主要均线下方；这不是把31.76元写成自动买点，也不是对后续价格的预测。\srcid{PR-02}

\section{成交、换手与流通口径：有数据，也不能偷换定义}

2026年7月22日，成交量为1,083,698手（1.0837亿股）、成交额36.04亿元、换手率1.44\%。近5/20/60个交易日平均换手率分别为1.38\%/1.67\%/1.57\%，平均成交额为34.14/43.85/44.08亿元。当日成交活跃度低于20日均值，但与5日均值接近，不能单独解释为资金转向。\srcid{PR-02}

\begin{exhibitbox}[流动性与股本：可验证字段及其边界]
\small
\begin{tabularx}{\exhibitboxwidth}{L{3.3cm}R{3.0cm}L{4.0cm}X}
\toprule
\textbf{字段（截止2026-07-22）} & \textbf{数值} & \textbf{交叉检查} & \textbf{允许的解释} \\
\midrule
流通A股本 / 总股本 & 75.2562 / 75.2562亿股 & 与公司Q1精确股本一致 & 当日换手率的供应商分母 \\
流通市值 / 总市值 & 2,484.96 / 2,484.96亿元 & 33.02元$\times$75.2562亿股 & 市场锚与市值口径 \\
成交量 / 成交额 / 换手 & 1.0837亿股 / 36.04亿元 / 1.44\% & 成交股数$\div$流通股本=1.4394\% & 活跃度，不识别资金类型 \\
前十大流通股东 & 41.7680亿股；占55.5011\% & 最新披露期为2026-03-31 & 集中度观察，不倒推可交易筹码 \\
严格自由流通股本 & 不发布 & 无统一、可复核的战略/锁定/长期持有分类 & 不以“流通股本”冒充自由流通股本 \\
\bottomrule
\end{tabularx}
\sourcenote{PR-02、CF-04；供应商流通字段可与量/额和价格/股本机械勾稽。2026-03-31的股东表不是2026-07-22实时持仓。}
\end{exhibitbox}

这里的关键不是再造一个貌似精确的“自由流通”数字。公开快照给出了可复算的流通A股本和换手率；但十大流通股东中混有国有平台、香港中央结算和ETF等不同性质账户，披露并未按即时可交易性统一分类。因此本报告发布流通A股本，不把它改称为严格自由流通股本，也不用“总股本减前十”倒推一个虚假的自由流通数。\srcid{PR-02}

\section{融资融券、沪股通与基金：区分当日流量和季度存量}

融资融券使用上交所明细和东方财富结构化序列交叉核验，最新可得交易日为2026年7月21日（7月22日收盘前的最后完整数据）。融资余额75.67亿元，单日融资净偿还4,933万元；融券余量38.43万股、融券余额1,266.5万元。融资余额约为当日市值的3.05\%，而融券余额在两融余额中很小。该组合反映杠杆多头存量仍大、短线去杠杆在发生；它不证明基本面已经转弱，也不能与基金持仓直接相加。\srcid{PR-03}

沪股通持股采用上交所可得的季度末持有量：2025年12月31日1.4334亿股，2026年3月31日1.7108亿股，2026年6月30日2.4093亿股，后者约占流通A股3.20\%。这说明季度维度的沪股通存量上升，但并非7月22日的实时北向净买入；针对2026年7月21日的日度端点已保存空响应，故不以季度存量冒充日度流。\srcid{PR-03}

基金持仓的最新可得报告期为2026年6月30日：162只基金合计持有1.7335亿股、对应市值58.58亿元，占流通A股2.30\%，较2026年3月31日的112只基金、1.5608亿股增加1,727.8万股（+11.07\%）。这是季末披露口径，受净申赎、估值和持仓披露范围影响，只可作为机构存量的滞后快照。\srcid{PR-03}

\begin{exhibitbox}[资本定位：日期、数值与不能推出的结论]
\small
\begin{tabularx}{\exhibitboxwidth}{L{2.4cm}L{2.1cm}L{4.1cm}X}
\toprule
\textbf{指标} & \textbf{日期} & \textbf{可引用数值} & \textbf{正确使用边界} \\
\midrule
融资融券 & 2026-07-21 & 融资余额75.67亿元；当日净偿还0.49亿元；融券余额0.013亿元 & 日度杠杆资金快照；不等于所有机构资金 \\
沪股通持股 & 2026-06-30 & 2.4093亿股，占流通A股约3.20\%；较2025-12-31增加0.9759亿股 & 季末存量；不能写成7月22日北向净买入 \\
基金持仓 & 2026-06-30 & 162只基金、1.7335亿股、58.58亿元、占流通A股2.30\% & 季末披露；不能识别主动/被动或当日调仓 \\
龙虎榜 & 近一年查询截至2026-07-22 & 无600150记录；全历史最近记录为2021-09-02 & 仅说明该公开查询窗口无上榜记录，不等于不存在大额交易 \\
\bottomrule
\end{tabularx}
\sourcenote{PR-03；融资融券为交易日数据，沪股通与基金为报告/结算日数据，禁止跨频率相减或相加。}
\end{exhibitbox}

\section{二级市场结论如何改变执行}

市场数据没有改变基本面估值：33.02元仍接近33.07元目标价并位于31.64--35.96元基础公允区间内。但它把原来的空泛“只用估值纪律”具体化为三层约束。第一，前复权1个月/3个月收益为-10.68\%/-8.58\%，且价格低于MA10至MA250，故不能把单日反弹升级成趋势反转；第二，31.76元是最近5/20/60日共同低点，只是复核经营与估值的价格观察位，不是脱离基本面即可执行的买点；第三，融资余额回落而季度沪股通、基金存量上升，资金信号存在不同频率与不同含义，不能被压缩成“主力进场”。\srcid{PR-02, PR-03}

\begin{exhibitbox}[价格、估值和验证事件的联合纪律]
\small
\begin{tabularx}{\exhibitboxwidth}{L{3.2cm}L{4.2cm}X}
\toprule
\textbf{状态} & \textbf{市场与估值解释} & \textbf{研究动作} \\
\midrule
31.76元附近且基本面失效未触发 & 接近近期区间下沿，但仍低于主要中期均线 & 先核对下跌是否由H1质量、交付或现金新信息驱动；仅提高复核优先级 \\
31.64--35.96元基础区间 & 当前33.02元处于区间内，目标上行约0.15\% & 已有仓位持有观察；新增资金等待利润、核心毛利和OCF三重验证 \\
高于35.96元而基本面未升级 & 价格开始支付超越Base的持续期；资金存量不足以替代盈利证据 & 不追价，要求正式中报和2027--2028持续期证据 \\
H1硬失效或现金恶化 & 任何技术位都不能覆盖盈利分母下修 & 先重跑模型与Bear概率，再讨论价格位置 \\
\bottomrule
\end{tabularx}
\sourcenote{PR-02、PR-03及冻结估值模型；本表是研究监测纪律，不构成证券交易指令。}
\end{exhibitbox}

因此，二级市场信息不是一个与基本面脱节的“技术面附录”。它目前支持\textbf{持有/等待验证}而非追价：价格尚未显示中期重回强势，杠杆资金在短期回撤中下降；同时，沪股通与基金的季度存量并未给出必须立即降级的证据。真正会改变动作的仍是正式中报中利润、核心毛利、经营现金流与交付节奏的联合验证。对严格自由流通股本和北向日度净流，报告保留可验证的边界，不以不可得字段编造精度。
